\emph{I missed this lecture. This lecture and the previous one covered pages 61 to 80 in Shmuel's handwritten notes. The notes I'm including here for this lecture are copied from those notes as closely to word for word as I could, and are not based on the lecture itself.}
\section{Remnants of Stars - Exotic Objects}
Stars on the main sequence towards the end of their lives will become either Red Giants (If their mass is below 8 solar masses) or Supergiants (If their mass is above 8 solar masses). Red Giants then become Planetary Nebulas and eventually White Dwarfs. Supergiants collapse and explode in a Supenova and then end up as either a Neutron Star or a Black Hole.
\subsection{White Dwarfs}
In the 19-20 centuries it was determined that the nearby star Sirius (the brightest star in the night sky) is actually a pair of two stars, Sirius A and B, which is a white dwarf. From Kepler's law it was determined that Sirius B's mass is approximately 1 solar mass, but because it had low luminosity and high temperature, it was concluded that its radius is approximately \(\qty{6000}{\kilo\meter}\). This means its average density is
\begin{equation}
    \overline{\rho} = \frac{M}{\frac{4\pi}{3}R^3} = \qty{2e6}{\gram\per\centi\meter^3}=\qty{2}{\ton\per\centi\meter^3}
\end{equation}
in other words, a teaspoon of a white dwarf is as heavy as an elephant. At the end of the star's life the core shrinks (because the fusion of hydrogen stops). Therefore both the density and temperature of the core grow, and as a result Helium and Carbon begin to fuse in the processes
\begin{align}
    \ce{3He^4 -> C^{12} + \gamma} && \ce{He^4 + C^{12} -> O^{16} + \gamma}
\end{align}
the first process being called ``triple alpha'' which releases \(\qty{7.275}{\mega\electronvolt}\). Eventually the Helium is used up which further compresses and heats the core, but eventually something stops the shrinking. Pauli's exclusion principle forbids electrons from coming too close together.
\subsubsection{Which density does this start at?}
When the distance between electrons become of the order of magnitude of \(\frac{1}{2}\lambda_{DB}\) the wavefunctions star to `feel' each other. In terms of density,
\begin{equation}
    \rho_q \simeq \frac{m_p}{\ins{\frac{1}{2}\lambda_{DB}}^3}=\frac{8m_p\ins{3m_e k_B T}^{3/2}}{h^3} = \qty{e4}{}T_8^{3/2}\qty{}{\gram\per\centi\meter^3}
\end{equation}
where \(T_8 \equiv \frac{T}{\qty{e8}{\kelvin}}\). Once \(\rho \gtrsim \rho_q\) the `magic' happens. For the core of the sun, \(T_8 = 0.015\) which gives \(\rho_q = \qty{640}{\gram\per\centi\meter^3}\) but since the sun's core's density is \(\qty{150}{\gram\per\centi\meter^3}\) it is very far from the quantum limit. But when a star's core shrinks it is possible to reach it.
\subsubsection{At The Quantum Limit}
Particles' energies are distributed either by the Bose-Einstein distribution (for photons, gluons, etc) or the Fermi-Dirac distribution (for elcetrons, protons, neutrons, etc). The FD distribution at temperatures \(k_B T \ll \epsilon_f\) it is approximately a step function where all energies \(\epsilon \leq \epsilon_f\) are populated and the rest are not. The fermi energy is
\begin{equation}
    \epsilon_f = \frac{\hbar^2}{2m_e}\ins{3\pi^2 n_e}^{2/3}
\end{equation}
\emph{The lecture notes here develop the fermi energy, I skip it as it's included in the prerequesite course Statistical Mechanics.} In conclusion, we'll find particles with high energy and momentum regardless of the temperature. They're being forced to be such by Pauli's exclusion principle.
\subsubsection{The Pressure}
Earlier in the course we developed the equation
\begin{equation}\label{p_pre_assumption}
    P = \frac{1}{3}\int\limits_0^\infty \dd p pv\ins{\diff{n}{p}}
\end{equation}
if we assume the particles are non-relativistic we can substitute the velocity
\begin{equation}\label{pe}
    P_e = \frac{1}{3m_e}\int\limits_0^\infty \dd p p^2 \ins{\diff{n_e}{p}}
\end{equation}
the number of electrons as a function of their momentum is
\begin{equation}
    n_e = \frac{8\pi}{3}\ins{\frac{p}{h}}^3
\end{equation}
so the density is
\begin{equation}
    \diff{n_e}{p} = \begin{cases}
        8\pi\frac{p^2}{h^3} & p \leq p_f \\
        0 & p > p_f
    \end{cases}
\end{equation}
inserting this back into~\ref{pe}
\begin{align}\label{pe2}
    P_e &= \frac{1}{3m_e}\int\limits_0^{p_f}\dd p p^2\cdot 8\pi\frac{p^2}{h^3}\\
    &= \frac{8\pi}{3m_e h^3}\int\limits_0^{p_f}\dd p p^4 = \frac{8\pi}{15m_e h^3}p_f^5
\end{align}
inserting \(p_f\)
\begin{equation}
    \boxed{P_e = \ins{\frac{3}{8\pi}}^{2/3}\frac{h^2}{5m_e}n_e^{5/3}}
\end{equation}
To find the expression in terms of density of mass we notice that
\begin{equation}
    \rho = n_+ \cdot Am_p
\end{equation}
and if we assume the material is fully ionized we can write \(n_e = n_+ \cdot Z\) which gives us
\begin{equation}
    \rho = n_+ Am_p = \frac{n_e}{Z}Am_p = m_p n_e \frac{A}{Z}
\end{equation}
Therefore
\begin{equation}
    \boxed{n_e = \frac{\rho}{m_p}\frac{Z}{A}}
\end{equation}
inserting this into \refeq{pe2}
\begin{equation}
    \boxed{P_e = \ins{\frac{3}{\pi}}^{2/3}\frac{h^2}{20m_e m_p^{5/3}}\ins{\frac{Z}{A}}^{5/3}\rho^{5/3}}
\end{equation}
For a typical White Dwarf, \(P_e \approx \qty{3e22}{\dyne\per\centi\meter^2}\) which is 5 orders of magnitude more than the sun's core.
\subsubsection{An Estimate For a Typical Radius}
Using the hydrostatic equation
\begin{equation}
    \diff{P}{r} = -\frac{GM_r}{r^2}\rho
\end{equation}
and approximating
\begin{align}
    P &\propto r^\beta \implies \diff{P}{r} \sim \frac{P}{r} \implies P\approx \frac{GM_\rho}{r} \\
    \rho &\approx \frac{M}{r^3} \implies \approx\frac{GM^2}{r^2}
\end{align}
Some development leads to
\begin{equation}
    \boxed{r_{WD} = \frac{h^2}{20 m_e m_p^{5/3}G}\ins{\frac{Z}{A}}^{5/3}M^{-1/3}}
\end{equation}
inserting constants
\begin{equation}
    \boxed{r_{WD} \approx \qty{e9}{\centi\meter} \cdot \ins{\frac{Z}{A}}\ins{\frac{M}{M_\mathSun}}^{-1/3}} 
\end{equation}
\subsubsection{What is \(P_e\) at the relativistic limit?}
Inserting \(v=c\) into \refeq{p_pre_assumption} we find
\begin{equation}
    \boxed{P_e = \ins{\frac{3}{8\pi}}^{1/3}\frac{hc}{4m_p^{4/3}}\ins{\frac{Z}{A}}^{4/3}\rho^{4/3}}
\end{equation}
to find the mass-radius relation in the relativistic case we'll equate the hydrostatic pressure to the degeneracy pressure and find
\begin{equation}
    \frac{GM^2}{r^4} \approx d\rho^{4/3}
\end{equation}
if we insert the full expression for the density we'll find that the radius cancels out! This means we can't have an expression to connect the radius and mass! We find only a condition on the mass:
\begin{equation}
    \boxed{M_{ch} = 0.21 \ins{\frac{hc}{Gm_p}}^{3/2}\ins{\frac{Z}{A}}^2 m_p}
\end{equation}
This is called the Chandrasekhar mass, and it is the maximum allowed mass for a white dwarf. If \(M>M_{ch}\) the degeneracy pressure doesn't hold it together and it will collapse further. It's interesting to note that the Chandrasekhar mass has constants from different fields of physics.
\subsubsection{Cooling of White Dwarfs}
When White Dwarfs are created they are very hot (\(T\sim\qty{e9}{\kelvin}\)), and they cool down through black body radiation. To find the time it takes to cool down we can equate the luminosity of the star to the time derivative of the ideal gas law
\begin{equation}
    -\sigma T^4\cdot4\pi r^2 = \frac{M}{m_p}k_B \diff{T}{t}
\end{equation}
Therefore
\begin{align}
    t_{cool}  &= \int\limits_0^{t_{cool}}\dd t\\
    &= \frac{\ins{M/m_p}k_B}{4\pi r^2\sigma}\int\limits_{T_0}^{T_f}\frac{\dd T}{T^4}\\
    &= -\frac{\ins{M/m_p}k_B}{4\pi r^2\sigma}\ins{-\frac{1}{3}}\ins{\frac{1}{T_f^3}-\frac{1}{T_0^3}}
\end{align}
setting the initial inverse temperature cubed to be negligible we find
\begin{equation}
    t_{cool} = \frac{\ins{M/m_p}k_B}{12\pi r^2\sigma}\frac{1}{T_f^3}
\end{equation}
so the time it would take a white dwarf to cool to 1000 kelvin would be
\begin{equation}
    t_{cool} = \qty{15}{\giga\year}\ins{\frac{M}{M_\mathSun}}\ins{\frac{4}{\qty{4000}{\kilo\meter}}}^{-2}
\end{equation}
which means it'd take longer than the age of the universe. In practive the cooling time is even longer. For this reason white dwarfs appear white (or teal).
\subsection{What Happens To More Massive Stars (\(M>8M_\mathSun\))?}
After the fusion of hydrogen ends, the pressure goes down, their core shrinks, it gets hotter, and fusion of the next element begins. First the triple alpha reaction, until it runs out too and the process happens again. The next stage is fusion of Carbon into Oxygen and Oxygen into Neon
\begin{align}
    \ce{He^4 + C^{12} -> O^{16} + \gamma} && \ce{He^4 + O^{16} -> Ne^{20} + \gamma} \\
    \ce{2C^{12} -> Ne^{20} + He^4 + \gamma} && \ce{2C^{12} -> Ne^{23} + p} && \ce{2C^{12} -> Mg^{23} + n}
\end{align}
and so on. Each stage is shorter than the last
\begin{table}[!h]
    \centering
    \begin{tabular}{l|c}
        H & \qty{5}{\mega\year}\\
        He & \qty{0.5}{\mega\year}\\
        C & \qty{500}{\year}\\
        Ne & \qty{1}{\year}\\
        Si & \qty{1}{} day
    \end{tabular}
\end{table}
this continues until it reaches Iron. From this point fusion no longer takes place, which destabilizes the core. There is no source of pressure which causes the core to collapse and heat up. This causes two processes to take place, one of which is called photodisintegration, which is the process
\begin{align}
    \ce{\gamma + Fe -> 13He + 4n} && \ce{\gamma + He -> 2p + n}
\end{align}
the first process consumes \(\qty{124}{\mega\electronvolt}\) and the second \(\qty{28}{\mega\electronvolt}\). If we assume a white dwarf converted all its mass to iron the number of iron atoms is
\begin{equation}
    N_{Fe} = \frac{M_{ch}}{m_{Fe}} \sim \qty{3e55}{}
\end{equation}
So the total energy to be lost is
\begin{equation}
    E_{loss} \sim \qty{2e52}{\erg}
\end{equation}
which is a massive amount of energy, equivalent to the total energy radiated by the sun through its entire life.